% teaching_statement.tex -- Teaching Statement (draft)
% Compile with: lualatex teaching_statement.tex
\documentclass[10pt]{article}

\usepackage[letterpaper,top=50pt,bottom=50pt,left=58pt,right=58pt]{geometry}
\usepackage{fontspec}
\usepackage[HTML]{xcolor}
\usepackage{hyperref}

\IfFontExistsTF{Calibri}
  {\setmainfont{Calibri}[Ligatures={CommonOff,TeXOff}]}
  {\setmainfont{Carlito}[Ligatures={CommonOff,TeXOff}]}

\definecolor{heading}{HTML}{365F91}
\definecolor{rule}{HTML}{444444}
\definecolor{tagline}{HTML}{333333}
\hypersetup{colorlinks=true, urlcolor=heading, linkcolor=heading}

% Font sizes (as in cv.tex, with slightly larger body text for continuous prose)
\newcommand{\namesize}{\fontsize{16pt}{19.5pt}\selectfont}
\newcommand{\headsize}{\fontsize{10.5pt}{12.8pt}\selectfont}
\newcommand{\bodysize}{\fontsize{10.5pt}{13.6pt}\selectfont}

% Section heading: as in cv.tex
\newcommand{\cvsection}[1]{%
  \par\vspace{8pt}%
  {\headsize\bfseries\color{heading}\strut#1\par}%
  \nopagebreak\vspace{1pt}%
  {\color{rule}\hrule height 0.75pt}%
  \nopagebreak\vspace{3pt}%
  \prevdepth=2.9pt\relax
}

\setlength{\parindent}{0pt}
\setlength{\parskip}{6pt}
\pagestyle{empty}

\begin{document}
\bodysize
\raggedright
\hyphenpenalty=10000
\frenchspacing
\clubpenalty=10000
\widowpenalty=10000

{\centering
  {\namesize\bfseries\color{heading}\strut Teaching Statement\par}\vspace{2pt}
  {\headsize\color{tagline}Michael Lehn, Ph.D. · Senior Lecturer, Institute of Numerical Mathematics, Ulm University\par}
}

% ===========================================================================
\cvsection{THE PROBLEM BEFORE THE SOLUTION}

One principle shapes almost everything I do as a teacher: I do not explain a solution to a problem that students do not yet understand. Before I introduce a method, a data structure, or a piece of hardware, I want students to see the problem it solves and why that problem matters. Once the problem is clear, the solution is no longer an arbitrary rule to be memorized; it becomes a natural next step that students can follow, question, and often anticipate. Teaching from first principles follows from this: if every new concept has to answer a question students already have, a course has to begin where the questions begin.

Since 2008 I have taught at Ulm University in three settings: required mathematics courses for science and engineering students, numerical mathematics, and computing courses, including a two-course sequence in high-performance computing that I designed.

% ===========================================================================
\cvsection{SEEING THE WHOLE MACHINE: HPC0}

\emph{Introduction to High Performance Computing} (HPC0), an elective, starts from a question that is easy to ask and hard to answer: what actually happens when a program runs? Bottom-up, students use the CircuitVerse simulator to build digital circuits from basic logic gates up to a small ALU of their own. They then explore processor design by developing and experimenting with an instruction set step by step, guided by videos and classroom exercises. For the rest of the course they use a provided instruction set with ulm-generator, a toolchain I developed, which supplies the processor components and generates an assembler, a simulator, and a debugger. Having built an ALU and experimented with an instruction set themselves, students can reason about the remaining components without constructing each one. Top-down, they program in ABC, a small C-like language I developed for teaching, and implement the parts of a compiler as regular assignments, from symbol tables and a lexer to a parser and code generators. Toward the end of the semester the two strands meet: the compiler generates code for this target architecture.

The same principle governs individual topics. Finite state machines are not introduced as a definition: students first encounter state-dependent behavior in the modal text editor they use every day, then need exactly this idea to write a lexer, and only then meet the formal concept. HPC0 is taught as a flipped classroom. Students prepare with short videos, and class time is spent on worksheets, while a teaching assistant and I help students develop, test, and understand their solutions. The worksheets, videos, and lecturer notes are \href{https://github.com/michael-lehn/not-abc/tree/main/hpc0-sessions}{openly available in English}.

% ===========================================================================
\cvsection{MEASURING BEFORE OPTIMIZING: HPC1}

\emph{High Performance Computing 1} is a graduate course that I designed and lead; part of the teaching is delegated to a co-instructor. Its central problem is the gap between what the hardware can do and what straightforward code achieves. Students measure that gap themselves by comparing their own implementations with the theoretical peak performance of the machine, and the cache hierarchy, blocking, and packing are introduced as answers to a problem they have already seen in their own numbers. Step by step, they build their own optimized matrix–matrix product, integrate and tune provided SIMD micro-kernels, and use their kernels in a blocked LU factorization. When it still lags behind a vendor library, the next step is to test a hypothesis about where the time is spent. Parallel programming follows the same pattern. In guided exercises, students analyze and repair deliberately simple thread-pool prototypes built on C++ threads, for example by explaining why all jobs end up on a single thread, then use the course library's thread pool to parallelize matrix computations before moving on to OpenMP, MPI, and CUDA (\href{https://www.mathematik.uni-ulm.de/numerik/hpc/ws25/}{session materials}).

% ===========================================================================
\cvsection{HARDWARE AND SOFTWARE TOGETHER}

\emph{Application-Oriented Software Development} is a required undergraduate lab for students of electrical and biomedical engineering. Students begin by typing a short machine-code program in Intel HEX format into a text editor, flashing it onto a microcontroller, and observing the signal on a pin with an oscilloscope. They then move to assembly, C, and the Arduino framework, still working in the terminal with Makefiles, and build controllers on breadboards with potentiometers and eventually an inertial measurement unit. The controllers drive games that the students write in Python using object-oriented design.

% ===========================================================================
\cvsection{MATHEMATICS FOR DIFFERENT AUDIENCES}

\emph{Higher Mathematics 1–3} is a required sequence for students of physics, engineering, and computational science, for whom mathematics is a tool rather than their major. Here I teach traditional lectures, and the problem-first principle means connecting each topic to the questions it answers, both within mathematics and in the students' own fields. In numerical mathematics, at every level, I want students to understand why a method is needed before they derive it. My \href{https://github.com/michael-lehn/numerical-linear-algebra-lecture-notes}{lecture notes for Numerical Linear Algebra} follow this pattern. Before any theory of floating-point arithmetic, they ask why one should not simply solve problems exactly. They answer with a small experiment: two formulas for $\sqrt{x+1}-\sqrt{x}$ that are identical in exact arithmetic, one of which retains only about six correct digits in double precision for $x=10^{10}$. Only then do the notes develop machine numbers and rounding-error analysis, and return to explain the result as cancellation. Later chapters use the same approach; for example, a simple $2\times 2$ matrix without an LU factorization motivates pivoting. I~also teach \emph{Mathematical Software}, a required programming lab for future secondary-school mathematics teachers.

% ===========================================================================
\cvsection{ONE COURSE, DIFFERENT STUDENTS}

HPC0 brings together bachelor's students in mathematics and computational science and master's students in physics, electrical engineering, and communications and computer engineering; in 2026, two secondary-school students from my extracurricular courses also completed it successfully. I cannot assume that these students start with the same knowledge: a shared problem gives them a common starting point, and in class I support each of them individually. In Application-Oriented Software Development I designed core tasks and optional extension tasks, so that students with different prior knowledge and working speeds can all take part meaningfully. In the lab, students work in pairs and have to agree on how their hardware and code fit together. When they propose competing explanations or solutions, I encourage them to test and discuss the alternatives rather than ask me which is right; taking a classmate's idea seriously enough to test it is a habit I hope they carry beyond computing.

% ===========================================================================
\cvsection{EVIDENCE AND REFLECTION}

The student evaluations available to me cover 30 offerings of my courses since 2008. The median overall rating is 5.2 on a six-point scale, with 6 as the best value, and most reports that include a comparison place my course above the average of mathematics lectures. Since recent courses had few respondents and not every offering produces a report, I treat these results as indications. The criticisms are at least as useful, and they are consistent: in large lectures, students ask for a slower pace and for complete written lecture notes, and in the computing courses they point to a heavy workload. Balancing ambitious projects against a realistic workload is what I think about most when I plan a course, because the projects students value most also cost them the most time.

% ===========================================================================
\cvsection{MENTORING AND SERVICE}

I have supervised about 45 bachelor's and master's theses since 2008, mostly in numerical mathematics and high-performance computing. Outside the university, I have mentored projects in Jugend forscht, Germany's national youth science competition, since 2008, and I co-direct the Ulm location of the Schülerforschungszentrum Südwürttemberg, a student research center where I teach extracurricular courses. Since 2002 I have regularly been elected as an academic staff representative to the study commission for the mathematics degree programs.

In every course, from large service lectures to computing courses built from the ground up, my aim is the same: students should understand the problem well enough that the solution makes sense before I present it.

\end{document}
